% !TeX root = ../../Book1.tex
% 纲要标题：### 19.2 Galois 对应
% 纲要位置：计划纲要.md，第 654 行。

\section{Galois 对应}
\label{b1:sec:1902}

% 纲要内容（仅作写作计划，尚非教材正文）：
%
% * 中间域与子群的反向对应
% * 次数、指数与群阶
% * 正规子群、正规子扩张与商群
% * 紧接对应定理，以 \(X^3-2\) 的分裂域枚举全部子群、不动域、次数、正规性和包含关系
%

现在固定有限 Galois 扩张 \(L/K\)，记 \(G=\operatorname{Gal}(L/K)\)。给定中间域 \(E\)，考察逐点固定 \(E\) 的自同构所组成的子群 \(\operatorname{Gal}(L/E)\)；给定子群，则取其共同不动元。两个构造反转包含关系；要证明它们互逆，还须使用不动域定理中的精确次数。

\subsection{中间域与子群的反向对应}
\label{b1:subsec:1902:01}

下述对应称为\term[Galois-duiying]{Galois 对应}[Galois correspondence]。
\begin{theorem}\label{b1:thm:finite-galois-correspondence}
有限 Galois 扩张 \(L/K\) 的中间域与 \(G=\operatorname{Gal}(L/K)\) 的子群，通过
\[
E\longmapsto\operatorname{Gal}(L/E),\qquad H\longmapsto L^H
\]
建立互逆的、反转包含关系的一一对应。对任意中间域 \(E\)，扩张 \(L/E\) 都是 Galois 扩张。
\end{theorem}
\begin{proof}
先说明最后一句。每个元素在 \(E\) 上的最小多项式整除它在 \(K\) 上的最小多项式，故可分；任意 \(E\) 嵌入也是 \(K\) 嵌入，所以由\cref{b1:thm:normal-embedding-criterion}，它保持 \(L\)，故 \(L/E\) 正规。于是 \(H=\operatorname{Gal}(L/E)\) 满足 \(|H|=[L:E]\)。不动域定理给出 \([L:L^H]=|H|\)，结合 \(E\subseteq L^H\) 以及塔式公式，得到 \(E=L^H\)。

反向给定任意 \(H\leq G\)，\cref{b1:thm:artin-fixed-field}直接给出 \(\operatorname{Gal}(L/L^H)=H\)。这证明两个构造互逆。若 \(E_1\subseteq E_2\)，固定较大域的条件更强，所以 \(\operatorname{Gal}(L/E_2)\subseteq\operatorname{Gal}(L/E_1)\)；子群越大，不动域越小。故对应反转包含。
\end{proof}
不能把最后一句反过来读成每个 \(E/K\) 都正规。例如三次多项式的非正规根域可以作为其 Galois 分裂域的中间域。

\subsection{次数、指数与群阶}
\label{b1:subsec:1902:02}

对对应的 \(E=L^H\)，不动域定理和塔式公式分别给出
\begin{equation}\label{b1:eq:galois-degree-index}
[L:E]=|H|,\qquad [E:K]=[G:H].
\end{equation}
若 \(E_1\subseteq E_2\) 对应 \(H_1\supseteq H_2\)，则 \([E_2:E_1]=[H_1:H_2]\)。这一等式计算的是群的指数，不要求相应子群正规。

\ProseParagraphBreak
例如 \(L=\mathbb F_{q^6}\)、\(K=\mathbb F_q\) 时，\cref{b1:thm:finite-field-automorphisms}给出 \(G=\langle\sigma\rangle\cong C_6\)。指数为 \(d\) 的子群 \(\langle\sigma^d\rangle\) 对应 \(\mathbb F_{q^d}\)，其中 \(d\mid6\)。这与\cref{b1:thm:finite-field-subfields}的直接计算一致；有限域的子域分类并不依赖本定理，因而这里是对两种方法的比较。

\subsection{正规子群、正规子扩张与商群}
\label{b1:subsec:1902:03}

\begin{theorem}\label{b1:thm:galois-quotient}
设 \(L/K\) 为有限 Galois 扩张，\(G=\operatorname{Gal}(L/K)\)，\(H\leq G\)，并令 \(E=L^H\)。则 \(E/K\) 为 Galois 扩张当且仅当 \(H\) 是 \(G\) 的正规子群。此时限制映射 \(G\rightarrow\operatorname{Gal}(E/K)\) 是满同态，且
\[
\operatorname{Gal}(E/K)\cong G/H.
\]
\end{theorem}
\begin{proof}
对任意 \(\sigma\in G\)，直接检验得到
\[
\sigma(E)=L^{\sigma H\sigma^{-1}}.
\]
事实上，若 \(x\in E\)，则 \(\sigma h\sigma^{-1}\) 固定 \(\sigma(x)\)；反向对 \(\sigma^{-1}\) 作同样检验。若 \(H\) 正规，则每个 \(\sigma\) 都保持 \(E\)。任意 \(E\) 的 \(K\) 嵌入能由\cref{b1:thm:embedding-extension}延拓到 \(L\) 的 \(K\) 嵌入，而 \(L/K\) 正规使此延拓属于 \(G\)。它保持 \(E\)，所以嵌入判据说明 \(E/K\) 正规；\(E/K\) 的可分性来自 \(L/K\)，于是为 Galois 扩张。

若 \(E/K\) 正规，每个 \(\sigma\in G\) 都保持 \(E\)，从而 \(L^{\sigma H\sigma^{-1}}=E=L^H\)。Galois 对应的单射性给出 \(\sigma H\sigma^{-1}=H\)，所以 \(H\) 正规。此时限制映射 \(G\rightarrow\operatorname{Gal}(E/K)\) 的核是恰好固定 \(E\) 的 \(H\)；任意目标自同构均能延拓到 \(L\)，故映射满射。\cref{b1:thm:first-isomorphism}遂把商群 \(G/H\) 同构到这个实际像。
\end{proof}
正规性确保限制后仍是 \(E\) 到自身的映射；若 \(E/K\) 不正规，某些 \(\sigma\) 会把它送到另一个共轭子域，不能对整个 \(G\) 定义到 \(\operatorname{Aut}_K(E)\) 的限制同态。

\ProseParagraphBreak
对上述 \(E=L^H\)，逐点固定 \(E\) 与将 \(E\) 保持为集合是不同的条件。记 \(N_G(H)=\{\,\sigma\in G\mid\sigma H\sigma^{-1}=H\,\}\) 为 \(H\) 在 \(G\) 中的正规化子。Galois 对应及证明中的共轭公式分别给出
\[
\{\,\sigma\in G\mid\sigma|_E=\operatorname{id}_E\,\}
=\operatorname{Gal}(L/E)=H,
\qquad
\{\,\sigma\in G\mid\sigma(E)=E\,\}=N_G(H).
\]
恰是 \(N_G(H)\) 中的自同构才能限制成 \(E\) 的 \(K\) 自同构；\cref{b1:ex:19-normalizer-automorphisms}将用它计算 \(\operatorname{Aut}_K(E)\)。

\subsection{中间域的枚举}
\label{b1:subsec:1904:01}

先以 \(X^3-2\) 的分裂域把对应定理用在一个非交换 Galois 群上：求出全部子群及其不动域，并核对次数、正规性和包含关系。

\begin{proposition}\label{b1:prop:cubic-splitting-field-correspondence}
令 \(a=\sqrt[3]2\in\mathbb R\)，\(\zeta\) 为本原三次单位根，并置 \(L=\mathbb Q(a,\zeta)\)。则 \([L:\mathbb Q]=6\)，且 \(G=\operatorname{Gal}(L/\mathbb Q)\cong S_3\)。除 \(\mathbb Q\) 与 \(L\) 外，全部中间域为唯一的二次域 \(\mathbb Q(\zeta)\) 和三个互异的三次域 \(\mathbb Q(\zeta^j a)\)、\(j=0,1,2\)。前者在 \(\mathbb Q\) 上正规，后三者均不正规。
\end{proposition}
\begin{proof}
由\cref{b1:thm:eisenstein}，\(X^3-2\) 在 \(\mathbb Q\) 上不可约，故 \([\mathbb Q(a):\mathbb Q]=3\)。域 \(\mathbb Q(a)\) 包含于 \(\mathbb R\)，而 \(\zeta\) 非实且满足 \(X^2+X+1=0\)，所以 \([L:\mathbb Q(a)]=2\)。这是一个特征零多项式的分裂域，因而为六次 Galois 扩张。群 \(G\) 忠实置换三个根 \(a,\zeta a,\zeta^2a\)，由阶数知它给出全部 \(S_3\)。

记循环置换三个根的自同构为 \(s\)，复共轭为 \(t\)，则
\[
s(a)=\zeta a,\quad s(\zeta)=\zeta,\qquad
t(a)=a,\quad t(\zeta)=\zeta^{-1}.
\]
它们满足 \(s^3=t^2=1\)、\(tst=s^{-1}\)，且生成 \(G\)。\(G\) 的全部子群为 \(G\)、\(\langle s\rangle\)、\(\langle s^{2j}t\rangle\)（\(j=0,1,2\)）及 \(1\)：真非平凡子群只能有二阶或三阶，分别由三个换位或一个三循环生成。

自同构 \(s\) 固定 \(\zeta\)，而 \(s^{2j}t\) 固定 \(\zeta^j a\)。由\eqref{b1:eq:galois-degree-index}，这两个子群的不动域在 \(\mathbb Q\) 上分别有次数二与三；\(\mathbb Q(\zeta)\) 与 \(\mathbb Q(\zeta^j a)\) 也分别有这些次数，故相应包含必为相等。对应的三个二阶子群互异，所以三个三次域也互异。子群 \(\langle s\rangle\) 正规，三个二阶子群彼此共轭而不正规，\cref{b1:thm:galois-quotient}遂给出各中间域的正规性。
\end{proof}

例如取 \(E=\mathbb Q(\zeta)\)。全部六个 \(G\) 中的自同构都将 \(E\) 保持为集合，但只有 \(\langle s\rangle\) 的三个元素逐点固定 \(E\)；复共轭 \(t\) 把 \(\zeta\) 送到 \(\zeta^{-1}\)，所以不逐点固定它。

\cref{b1:tab:cubic-galois-correspondence}集中列出上述对应。对每一行都有 \([E:\mathbb Q]=[G:H]\)、\([L:E]=|H|\)；最后一列同时说明 \(H\) 是否为 \(G\) 的正规子群，因为这里的全部扩张均可分。
\begin{table}[htbp]
\centering
\caption[\(X^3-2\) 分裂域的 Galois 对应]{\(X^3-2\) 分裂域的 Galois 对应。其中 \(j=0,1,2\)。}
\label{b1:tab:cubic-galois-correspondence}
\begin{tabular}{@{}llccl@{}}
\toprule
子群 \(H\) & 不动域 \(E=L^H\) & \([E:\mathbb Q]\) & \([L:E]\) & \(E/\mathbb Q\) 正规 \\
\midrule
\(G\cong S_3\) & \(\mathbb Q\) & \(1\) & \(6\) & 是 \\
\(\langle s\rangle\cong C_3\) & \(\mathbb Q(\zeta)\) & \(2\) & \(3\) & 是 \\
\(\langle s^{2j}t\rangle\cong C_2\) & \(\mathbb Q(\zeta^j a)\) & \(3\) & \(2\) & 否 \\
\(1\) & \(L\) & \(6\) & \(1\) & 是 \\
\bottomrule
\end{tabular}
\end{table}

\cref{b1:fig:cubic-intermediate-fields}将表中的三个三次域分别画出；图中较低的域包含于较高的域，边上的数字是相邻两域的扩张次数。
\begin{figure}[htbp]
\centering
\begin{tikzpicture}[x=1cm,y=1cm,every node/.style={font=\small}]
\node (L) at (0,1.8) {\(L\)};
\node (Z) at (-5.5,0) {\(\mathbb Q(\zeta)\)};
\node (A) at (-1.8,0) {\(\mathbb Q(a)\)};
\node (B) at (1.8,0) {\(\mathbb Q(\zeta a)\)};
\node (C) at (5.5,0) {\(\mathbb Q(\zeta^2a)\)};
\node (Q) at (0,-1.8) {\(\mathbb Q\)};
\draw (L) -- node[pos=.55,fill=white,inner sep=1pt] {\scriptsize 3} (Z);
\draw (L) -- node[pos=.55,fill=white,inner sep=1pt] {\scriptsize 2} (A);
\draw (L) -- node[pos=.55,fill=white,inner sep=1pt] {\scriptsize 2} (B);
\draw (L) -- node[pos=.55,fill=white,inner sep=1pt] {\scriptsize 2} (C);
\draw (Z) -- node[pos=.45,fill=white,inner sep=1pt] {\scriptsize 2} (Q);
\draw (A) -- node[pos=.45,fill=white,inner sep=1pt] {\scriptsize 3} (Q);
\draw (B) -- node[pos=.45,fill=white,inner sep=1pt] {\scriptsize 3} (Q);
\draw (C) -- node[pos=.45,fill=white,inner sep=1pt] {\scriptsize 3} (Q);
\end{tikzpicture}
\caption[\(X^3-2\) 分裂域的中间域格]{\(X^3-2\) 分裂域的中间域格。边上的数字为较高的域在较低的域上的次数。}
\label{b1:fig:cubic-intermediate-fields}
\end{figure}

图中的四个非平凡真中间域两两互不包含：二次与三次之间由次数的整除性排除包含，三个不同的三次域之间也不能包含。把每条包含关系反向，就得到 \(G\) 的子群格。例如 \(\mathbb Q\subset\mathbb Q(\zeta)\subset L\) 对应 \(G\supset\langle s\rangle\supset1\)。这个枚举穷尽所有中间域，因为已经列出了全部子群；一般地，有限 Galois 扩张的中间域数必有限。
